% !PRINTED 205
\weberpage{205}
%<<CONT
\[
\psi_1 = a_1'\xi^2 - a_3'\eta^2
\]
\[
= a_0\,\frac{(\beta-\gamma)(\beta-\delta)(x-\alpha)^2 - (\alpha-\gamma)(\alpha-\delta)(x-\beta)^2}{\alpha-\beta}.
\]
In this expression, however, $\alpha - \beta$ can be cancelled from numerator and denominator, and
$\psi_1$ can be represented in two forms; to obtain $\psi_2$ and $\psi_3$, one then
need only permute $\beta$, $\gamma$, $\delta$ cyclically.

One thus finds
\[
\begin{aligned}
\frac{\psi_1}{a_0} &= (\gamma-\beta)(x-\alpha)(x-\delta) - (\alpha-\delta)(x-\beta)(x-\gamma) \\
&= (\delta-\beta)(x-\alpha)(x-\gamma) - (\alpha-\gamma)(x-\beta)(x-\delta)
\end{aligned}
\]
\weq{6}{
\begin{aligned}
\frac{\psi_2}{a_0} &= (\delta-\gamma)(x-\alpha)(x-\beta) - (\alpha-\beta)(x-\gamma)(x-\delta) \\
&= (\beta-\gamma)(x-\alpha)(x-\delta) - (\alpha-\delta)(x-\gamma)(x-\beta) \\
\frac{\psi_3}{a_0} &= (\beta-\delta)(x-\alpha)(x-\gamma) - (\alpha-\gamma)(x-\delta)(x-\beta) \\
&= (\gamma-\delta)(x-\alpha)(x-\beta) - (\alpha-\beta)(x-\delta)(x-\gamma).
\end{aligned}
}
The quantities $\psi_1^2$, $\psi_2^2$, $\psi_3^2$ may be regarded as the roots
of a cubic equation whose coefficients are covariants. One obtains
this equation easily from (4) in the form
\[
\left(z + \frac{H+4a_2'f}{3}\right)\left(z + \frac{H+4a_2''f}{3}\right)\left(z + \frac{H+4a_2'''f}{3}\right) = 0.
\]
By \S.\,65, (12), and making use of the expression (5) for $T^2$, this equation yields
\weq{7}{
z^3 + Hz^2 + \tfrac{1}{3}(H^2 - 16Af^2)z - T^2 = 0.
}
It is also of interest to form the discriminant of this equation. One
obtains it most simply from the expression
\[
\varDelta = (\psi_1^2 - \psi_2^2)^2(\psi_1^2 - \psi_3^2)^2(\psi_2^2 - \psi_3^2)^2,
\]
if one substitutes the expressions (4) and then, for $(a_2'-a_2'')^2(a_2'-a_2''')^2(a_2''-a_2''')^2$,
puts the discriminant of the equation \S.\,65, (12)
\[
27(4A^3 - B^2) = 3^6 D
\]
where then $D$ is the discriminant of $f$. One thus obtains
\weq{8}{
\varDelta = 2^{12} D f^6.
}
